\chapter{Generate the Catalog API}

The catalog service will be a small Rails API rather than a full browser-facing
application. Before generating it, we will preview the generator and create a
dedicated Unix account. Keeping the application files out of \texttt{root}'s
ownership makes later dependency and server commands safer and more
predictable.

\section{Exercise 6A: Preview the generator}

From the catalog jail's root console, ask Rails to show what it would generate:

\begin{lstlisting}[language=sh]
rails new /usr/local/www/catalog \
  --api \
  --minimal \
  --database=postgresql \
  --skip-bundle \
  --skip-git \
  --skip-docker \
  --skip-kamal \
  --skip-solid \
  --pretend
\end{lstlisting}

The \texttt{--pretend} flag runs the generator without writing anything. Check
that the destination still does not exist:

\begin{lstlisting}[language=sh]
test ! -e /usr/local/www/catalog && \
  echo "preview made no files"
\end{lstlisting}

The check should print \texttt{preview made no files}. The plain
\texttt{rails} command selects the jail's installed Rails version; naming a
specific patch version in every command would add noise without helping this
integration exercise.

\section{Exercise 6B: Create the application account}

Create a local group and user for the catalog service, followed by an empty
application directory owned by that account:

\begin{lstlisting}[language=sh]
pw groupadd catalog
pw useradd catalog -g catalog -m -s /bin/sh \
  -c "Catalog Rails application"
install -d -o catalog -g catalog -m 0755 \
  /usr/local/www/catalog
\end{lstlisting}

Verify both identities and the directory boundary:

\begin{lstlisting}[language=sh]
id catalog
stat -f '%Su:%Sg %Sp %N' /usr/local/www/catalog
\end{lstlisting}

The first command should identify the \texttt{catalog} user and group. The
second should begin with \texttt{catalog:catalog}, confirming that the service
account---not \texttt{root}---will own the generated application.

\section{Exercise 6C: Generate the application as its owner}

The preview established what Rails intends to create. Now run the same
generator as the service account, using \texttt{.} because the command starts
inside the application directory:

\begin{lstlisting}[language=sh]
su -l catalog -c 'cd /usr/local/www/catalog && \
  rails new . \
    --api \
    --minimal \
    --database=postgresql \
    --skip-bundle \
    --skip-git \
    --skip-docker \
    --skip-kamal \
    --skip-solid'
\end{lstlisting}

On FreeBSD, placing \texttt{-c} after the login name passes it to the user's
shell. The login shell therefore executes the quoted command as
\texttt{catalog}, rather than generating root-owned application files.

The \texttt{--skip-bundle} flag deliberately separates source generation from
dependency installation. Before Bundler changes anything, verify the generated
boundary and inspect the two central database dependencies:

\begin{lstlisting}[language=sh]
stat -f '%Su:%Sg %Sp %N' \
  /usr/local/www/catalog/Gemfile
grep -E '^gem "(rails|pg)"' \
  /usr/local/www/catalog/Gemfile
\end{lstlisting}

The \texttt{Gemfile} should still belong to \texttt{catalog:catalog}. Its Rails
entry should select the installed Rails 8.1 line, and its PostgreSQL entry
should use the \texttt{pg} adapter.

\section{Exercise 6D: Complete the bundle}

Install the two remaining direct dependencies from FreeBSD, then create a
lockfile from the gems already present in this dedicated application jail:

\begin{lstlisting}[language=sh]
pkg install -y rubygem-puma rubygem-debug
su -l catalog -c \
  'cd /usr/local/www/catalog && bundle lock --local'
su -l catalog -c \
  'cd /usr/local/www/catalog && bundle check'
\end{lstlisting}

The exact patch versions are not important to this tutorial. The final command
only needs to report that the bundle is satisfied. From here onward, the focus
returns to the service boundary: configuring this application to use the
PostgreSQL jail and exposing an API endpoint for the end-to-end test.
